% !PRINTED 207
\weberpage{207}

%<<CONT
setzt. Die Formel (2) des § 21 ergibt dann:

\weq{6}{C = (-1)^{(g_1+1)(g_2+1)\frac{n-1}{2}} n \frac{\vartheta_{g_1,g_2}(0,\omega)}{\vartheta_{g_1,g_2}(2\nu'\omega, n\omega)} \frac{\vartheta_{11}'(0, n\omega)}{\vartheta_{11}'(0, \omega)}.}

Die \gesperrt{Abelschen} Relationen erhält man einfach, indem man
in (5) $nu = 2\nu'\omega$ setzt, wodurch die rechte Seite verschwindet:

\weq{7}{\sum^{\nu} e^{\frac{2\nu\nu'\pi i}{n}} \frac{\vartheta_{11}\left(\frac{2\nu'\omega + 2\nu}{n}\right)}{\vartheta_{g_1,g_2}\left(\frac{2\nu'\omega + 2\nu}{n}\right)} = 0.}
% [sic] The source prints 2 in this exponential, although direct substitution
% in (5), and the downstream formulae (8)--(11), retain the factor 8.

Wendet man auf die hier vorkommenden $\vartheta$-Funktionen eine
lineare Transformation

\[
\begin{pmatrix} \alpha, & \beta \\ \gamma, & \delta \end{pmatrix}, \qquad \alpha\delta - \beta\gamma = 1
\]

an, so erhält man nach § 39 die allgemeine Formel:

\weq{8}{\sum^{\nu} e^{\frac{8\nu\nu'\pi i}{n}} \frac{\vartheta_{11}\left(\frac{2(\nu\alpha+\nu'\gamma) + 2(\nu\beta+\nu'\delta)\omega}{n}\right)}{\vartheta_{g_1,g_2}\left(\frac{2(\nu\alpha+\nu'\gamma) + 2(\nu\beta+\nu'\delta)\omega}{n}\right)} = 0.}

Diese Gleichungen gehen nun, wenn man $g_1, g_2 = 0, 1$ setzt
und die Bezeichnungen des § 42 einführt, unmittelbar in Relationen zwischen den Wurzeln $x_{\nu,\nu'}$ (§ 61) über:

\weq{9}{\sum^{\nu} e^{\frac{8\nu\nu'\pi i}{n}} x_{\nu\alpha+\nu'\gamma,\ \nu\beta+\nu'\delta} = 0,}

und den beiden speziellen Fällen:

\[
\begin{pmatrix} \alpha, & \beta \\ \gamma, & \delta \end{pmatrix} = \begin{pmatrix} 1, & 0 \\ 0, & 1 \end{pmatrix}, \qquad \begin{pmatrix} 0, & 1 \\ -1, & 0 \end{pmatrix}
\]

entsprechend:

\weq{10}{\sum^{\nu} e^{\frac{8\nu\nu'\pi i}{n}} x_{\nu,\nu'} = 0,}

\weq{11}{\sum^{\nu'} e^{-\frac{8\nu\nu'\pi i}{n}} x_{\nu,\nu'} = 0.}

Nach (5), (6) kann man auch den Wert dieser Summen bestimmen, wenn darin die $n^{\text{te}}$ Einheitswurzel

\[
e^{\frac{8\pi i}{n}}
\]
%CONT>>
